\section{Hipótesis}

Las hipótesis se formulan como afirmaciones falsables sobre la relación entre decisión distribuida, factibilidad física y coste operativo. Cada una se evalúa con mundos compartidos, métricas comparables y una lectura explícita del alcance. La tesis no exige que un único método domine todas las métricas; exige demostrar dónde aparece la mejora, dónde falla y qué parte del fenómeno queda fuera de la validación.

\begin{table}[H]
\centering
\footnotesize
\begin{tabularx}{\textwidth}{@{}p{0.10\textwidth}X X@{}}
\toprule
ID & Hipótesis & Criterio de lectura \\
\midrule
H1 & El reclutamiento distribuido puede cerrar coaliciones útiles para cargas heterogéneas sin asignación central completa. & Mayor satisfacción de demanda o menor brecha frente a referencias bajo los mismos mundos. \\
H2 & La capacidad efectiva cambia el problema respecto a la cardinalidad pura. & Una coalición con suficientes robots puede seguir siendo insuficiente si aporta poca capacidad mecánica. \\
H3 & La factibilidad \emph{wrench} evita falsos positivos físicos. & El residual de fuerza y torque debe separar coberturas escalares aparentes de contactos físicamente útiles. \\
H4 & La llegada segura al objetivo introduce un compromiso entre rapidez, energía y colisiones. & No se declara ganador único cuando la mejora en llegada aumenta coste o reduce margen de seguridad. \\
H5 & El transporte cooperativo exige evaluar la carga como cuerpo extendido. & La pose final, la integridad de formación y los despejes son tan relevantes como la asignación inicial. \\
H6 & La recuperación ante fallos requiere reparar coaliciones durante la ejecución, no solo reasignar al inicio. & La lectura combina éxito de recuperación, cargas perdidas, tiempo de reparación y seguridad posterior al evento. \\
H7 & La comunicación local y el sensado degradado pueden romper el transporte aunque la asignación inicial sea válida. & Se distinguen conectividad directa, conectividad temporal, retransmisión y éxito de transporte. \\
H8 & A escala de almacén, las soluciones centralizadas exactas pierden practicidad antes que las familias distribuidas o jerárquicas. & El criterio combina tiempo de ejecución, finalización, factibilidad \emph{wrench}, mensajes y coste por recurso. \\
\bottomrule
\end{tabularx}
\caption{Hipótesis principales de la memoria y criterio de lectura.}
\label{tab:hipotesis-principales}
\end{table}

\subsection*{Escalera modular de validación}
\label{sec:escalera-modular}

La memoria avanza por una escalera de dificultad creciente. Esta organización evita mezclar evidencias de distinta madurez: algunas conclusiones se apoyan en Monte Carlo de simulación, otras en pruebas algebraicas o en inspección visual, y otras quedan delimitadas como extensión.

\begin{table}[H]
\centering
\footnotesize
\begin{tabularx}{\textwidth}{@{}p{0.12\textwidth}p{0.28\textwidth}Y Y@{}}
\toprule
Módulo & Pregunta técnica & Evidencia esperada & Estado de lectura \\
\midrule
M0 & ¿Qué ocurre con reglas clásicas, centralizadas o puramente locales? & Línea base de coste, calidad y fallo bajo comunicación local. & Comparativo. \\
M1 & ¿Cómo se cierran quórums enteros para cargas con demanda mínima? & Reducción de coaliciones parciales y mejora de completitud. & Implementado y evaluado. \\
M2 & ¿Cómo cambia el problema cuando los robots tienen capacidades efectivas distintas? & Satisfacción de capacidad, coste físico y brecha frente a referencia. & Implementado y evaluado. \\
M3 & ¿La coalición genera el \emph{wrench} que exige la carga? & Residual, falsos positivos escalares y roles de contacto. & Implementado en simulación planar. \\
M4 & ¿Los robots llegan y mantienen seguridad durante el movimiento? & Tiempo, energía, colisiones, despeje y fallo por escenario. & Implementado y evaluado. \\
M5 & ¿La carga se mueve como cuerpo extendido hasta una pose final? & Éxito de transporte, formación, orientación y obstáculos. & Implementado y evaluado. \\
M6 & ¿El sistema repara fallos, batería y cargas inviables durante la misión? & Recuperación, cargas perdidas, detección y coste de reasignación. & Implementado y evaluado. \\
M7--M8 & ¿Qué ocurre con comunicación degradada y escala de almacén? & Conectividad temporal, sensado, rendimiento y frontera de intractabilidad. & Implementado en modelo mesoscópico. \\
\bottomrule
\end{tabularx}
\caption{Escalera modular usada para ordenar la validación.}
\label{tab:hipotesis-modulos}
\end{table}